\documentclass[10pt,a4paper]{amsart}
\usepackage[margin=19mm]{geometry}
\usepackage{amsmath,amssymb,mathtools}
\usepackage[scr=rsfs,cal=euler]{mathalfa}
\usepackage{xfrac}
\usepackage[unicode,hidelinks,bookmarks=false]{hyperref}
\newcommand{\ie}{i.e.\ }
\newcommand{\cf}{cf.\ }
\newcommand{\resp}{respectively\ }
\newcommand{\Subsection}{\subsection}
\providecommand{\nobreakdash}{\nobreak\hbox{-}}
\setlength{\emergencystretch}{2em}
\multlinegap0pt
\usepackage{amssymb}
\usepackage[shortlabels]{enumitem}
\usepackage{xcolor}
\usepackage[normalem]{ulem}
\usepackage{tikz}
\usepackage{subfig}
\usepackage[all]{xy}
\usepackage{mathtools}
\newtheorem{lem}{Lemma}
\usepackage{cleveref}
\crefname{lem}{Lemma}{Lemmas}
\def\bC{\mathbb{C}}
\def\bD{\mathbb{D}}
\def\bZ{\mathbb{Z}}
\renewcommand{\bar}[1]{\ol{#1}}
  \newcommand\bfM{{\bf M}}
\newcommand{\bs}{\backslash}
\newcommand{\ol}{\overline}	
   \def\trans{\mathrm{tr}}
\title{Baily--Borel compactifications of period images and the b-semiampleness conjecture}
\author{Benjamin Bakker, Stefano Filipazzi, Mirko Mauri, Jacob Tsimerman}
\date{Source: arXiv:2508.19215v2 (CC BY 4.0)}
\begin{document}
\maketitle

\noindent\emph{Source and licence.} Baily--Borel compactifications of period images and the b-semiampleness conjecture. Version arXiv:2508.19215v2 (CC BY 4.0), distributed by arXiv under the Creative Commons Attribution 4.0 International licence, http://creativecommons.org/licenses/by/4.0/, as recorded in the arXiv licence metadata for this exact version and shown on the abstract page https://arxiv.org/abs/2508.19215v2. Full licence notice: \texttt{licenses/arxiv-2508.19215-CC-BY-4.0.txt}.\par\smallskip

\noindent\textit{Editorial scope.} Counted unit: the article's lem stuff (source0.tex lines 1270--1277). The article's complete original proof of it is delivered with it (source0.tex lines 1278--1286, one \texttt{proof} environment of 9 lines). No mathematics is written, repaired or re-derived: the statement and the proof are the article's own lines, in the article's own notation, and no retained line is substituted or re-flowed.  Retained uncounted as the source-local context the proof needs: lines 1158--1159; lines 1244--1250.  Nothing was omitted from any retained span, so the package carries no omission record.  No retained line contains a citation, so the wrapper carries no bibliography and no bibliography entry.  No retained line contains a figure environment, an image-inclusion command or drawing code, so no image is delivered and the package declares no asset dependency.  Editorial formatting only, in addition: an AMS \texttt{amsart} preamble that restates the article's own declarations and macros used by the retained lines, namely the environments \texttt{lem} on separate counters, together with the standard packages those lines need.  The retained environments print in this document as Lemma 1 in order of appearance, whereas the article prints the same items with the numbering of its own full document.  The complete native source of the article as distributed in the arXiv e-print for this exact version is bundled unchanged as \texttt{sources/183941/source0.tex}.
    \begin{lem}\label{lem:unpol orbit}For any $x\in \bfM(\bZ)\backslash \bD$, there are finitely many points $x'\in \bfM(\bZ)\backslash \bD$ for which the associated integral Hodge structures $V(x),V(x')$ are isomorphic.
\end{lem}

\begin{lem}\label{lem:equiv reln}Let $X$ be a proper algebraic space and $R\subset X(\bC)\times X(\bC)$ a closed reflexive symmetric constructible relation.  Then there is a closed constructible subset $\Delta\subset X(\bC)$ such that: 
\begin{enumerate}
\item Every $x\in X(\bC)$ with non-constructible $R^e$-equivalence class is contained in $\Delta$.
\item The set of points $x\in \Delta$ with constructible $R^e$-equivalence class is contained in a countable union of nowhere dense constructible subsets of $\Delta$.
\end{enumerate}

\end{lem}

\begin{lem}\label{lem:equiv reln stuff}\hspace{.5in}
\begin{enumerate}
\item\label{lem:equiv reln stuff p0} For any connected constructible subset $Z$ of a $R_\trans$-equivalence class, the closure $\bar Z$ is contained in an $R_\mathrm{curve}$-equivalence class.  
\item\label{lem:equiv reln stuff p1} If a $R_\trans$-equivalence class is constructible, then it is closed and its connected components are $R_\mathrm{curve}$-equivalence classes.  
\item\label{lem:equiv reln stuff p2} For any point $x\in X(\bC)$ for which $V^\trans(x)$ is maximal rank, the $R_\trans$-equivalence class of $x$ is constructible.
\item\label{lem:equiv reln stuff p3} For any closed reflexive symmetric constructible relation $R\subset X(\bC)\times X(\bC)$ for which $R_\mathrm{curve}\subset R^e\subset R_\trans$, $R^e$ is a closed constructible equivalence relation whose equivalence classes are finite unions of $R_\mathrm{curve}$-equivalence classes.
\end{enumerate}
\end{lem}

\begin{proof}
For \eqref{lem:equiv reln stuff p0}, if a $R_\trans$-equivalence class contains a constructible set $Z$, then for any proper connected curve $C\subset X$ whose generic point is contained in $Z$, $V^\trans(C)$ as above is isotrivial, so $C$ is contained in a $R_\mathrm{curve}$-equivalence class, hence $C\subset Z$.  Thus $\bar Z$ is contained in a $R_\mathrm{curve}$-equivalence class.  Part \eqref{lem:equiv reln stuff p1} follows immediately from \eqref{lem:equiv reln stuff p0}.



For \eqref{lem:equiv reln stuff p2}, let $x\in X(\bC)$ be such a point and $Z$ its $R_\trans$-equivalence class.  Then for any $z\in Z\cap X_\Sigma$, $V^\trans(x)=V^\trans_\Sigma$.  Applying \Cref{lem:unpol orbit}, $Z\cap X_\Sigma$ is constructible, hence $Z$ is.

For \eqref{lem:equiv reln stuff p3}, by \Cref{lem:equiv reln} there is a closed constructible $\Delta\subset X(\bC)$ such that outside a countable union of nowhere dense constructible subsets $\Xi\subset \Delta$, each $R$-equivalence class is non-constructible. However applying \eqref{lem:equiv reln stuff p2} to (a resolution of) each irreducible component $\Delta_0$ of $\Delta$, each $R_\trans|_{\Delta_0}$-equivalence class in $\Delta_0$ of maximal rank transcendental part is constructible, as therefore is any $R_\trans|_\Delta$-equivalence class in $\Delta$ of maximal rank transcendental part.  If $\Delta$ were nonempty, there would then be a point $x\in\Delta\bs \Xi$ of maximal rank transcendental part, hence constructible $R_\trans|_\Delta$-equivalence class $Z$.  By \eqref{lem:equiv reln stuff p1}, $Z$ is partitioned into finitely many constructible $R_\mathrm{curve}$-equivalence classes, hence also into finitely many constructible $R^e$-equivalence classes.  This is a contradiction, so $\Delta=\varnothing$. 
\end{proof}

\end{document}
